\clearpage
\refstepcounter{section}
\addcontentsline{toc}{section}{\protect\numberline{}Scheda riassuntiva e risoluzione dei problemi}
\markboth{Scheda riassuntiva}{Scheda riassuntiva}
\begin{tikzpicture}
  \fill[navy,rounded corners=8pt] (0,0) rectangle (\textwidth,1.9cm);
  \node[anchor=west,font=\headfont\bfseries\fontsize{19}{23}\selectfont,text=white] at (.6cm,1.15cm) {Scheda riassuntiva dei comandi};
  \node[anchor=west,font=\headlight\small,text=white!80!teal] at (.6cm,.45cm) {Da tenere accanto alla tastiera durante il laboratorio};
\end{tikzpicture}
\vspace{4pt}

\begin{tcbraster}[raster columns=2,raster equal height=rows,raster column skip=10pt,raster row skip=8pt]
\begin{tcolorbox}[base,colback=teal!5,colbacktitle=teal,coltitle=white,title={Sull'host (\$)},top=13pt,fontupper=\small,unbreakable]
\renewcommand{\arraystretch}{1.15}
\begin{tabular}{@{}>{\leavevmode\color{navy}\ttfamily\footnotesize}p{3.75cm} >{\raggedright\arraybackslash\footnotesize}p{2.9cm}@{}}
kathara lstart & avvia il lab\\
kathara lclean & spegne il lab\\
kathara lrestart & riparte da zero\\
kathara connect pc1 & shell su pc1\\
kathara exec r1 "ip r" & un comando su r1\\
lconfig -n D --add LAN & collega D a una LAN\\
lconfig -n D --rm LAN & scollega D\\
kathara linfo [-t] [-w] & stato / topologia\\
kathara list & dispositivi attivi\\
kathara wipe & spegne tutto\\
kathara check & diagnostica\\
\end{tabular}
\end{tcolorbox}
\begin{tcolorbox}[base,colback=navy!5,colbacktitle=navy,coltitle=white,title={Nei dispositivi (\#)},top=13pt,fontupper=\small,unbreakable]
\renewcommand{\arraystretch}{1.15}
\begin{tabular}{@{}>{\leavevmode\color{navy}\ttfamily\footnotesize}p{4.15cm} >{\raggedright\arraybackslash\footnotesize}p{2.5cm}@{}}
ip a  /  ip -br a & indirizzi\\
ip a add IP/n dev eth0 & assegna IP\\
ip a del IP/n dev eth0 & toglie IP\\
ip link set eth0 up|down & accende/spegne\\
ip r  /  routel & tabella di routing\\
ip r add default via GW & default gateway\\
ip r add NET/n via NH dev IF & rotta statica\\
ip r del NET/n & toglie rotta\\
ip r get IP & quale rotta usa?\\
arp -n  /  ip neigh & cache ARP\\
ping -c N IP & raggiungibilità\\
traceroute [-n] [-z 1] IP & router sul percorso\\
\end{tabular}
\end{tcolorbox}
\end{tcbraster}

\subsection*{\color{navy}Il piano del lab in una riga}
\begin{center}
\topologia{.72}{0}{1}
\end{center}

\subsection*{\color{navy}Il ping non funziona: cosa controllare}
\begin{center}
\begin{tikzpicture}[
  st/.style={draw=navy,fill=navy!5,rounded corners=4pt,font=\headfont\small,align=left,text width=6.2cm,inner sep=5pt},
  fx/.style={draw=rosso,fill=rosso!6,rounded corners=4pt,font=\headfont\footnotesize,align=left,text width=6.4cm,inner sep=5pt},
  ar/.style={-{Stealth[length=6pt]},line width=1pt,color=grigio},
  n/.style={circle,fill=teal,text=white,font=\headfont\bfseries\small,inner sep=2pt,minimum size=.55cm}]
  \foreach \i/\q/\f in {
    1/{L'interfaccia ha l'indirizzo giusto ed è UP?}/{\texttt{ip -br a}; se manca: \texttt{ip a add \dots}, \texttt{ip link set eth0 up}},
    2/{Il mittente ha una rotta per la destinazione?}/{\texttt{ip r get IP}; \texttt{Network is unreachable} = manca il default gateway},
    3/{L'ARP verso il next hop riesce?}/{\texttt{arp -n}: voce \texttt{incomplete} = vicino spento, IP sbagliato o LAN sbagliata nel \texttt{lab.conf}},
    4/{I router hanno la rotta di andata?}/{\texttt{routel} su ogni router; messaggi \texttt{Net Unreachable} indicano quale router manca di rotta},
    5/{E la rotta di ritorno?}/{\texttt{100\% packet loss} senza errori: spesso manca la rotta per la rete del mittente (esercizio 4)},
    6/{Cosa passa davvero sul filo?}/{\texttt{kathara lconfig -n wireshark --add LAN} e cattura: si vede dove i pacchetti si fermano}}{
    \node[n] (n\i) at (0,-\i*1.32) {\i};
    \node[st,anchor=west] (q\i) at (.5,-\i*1.32) {\q};
    \node[fx,anchor=west] (f\i) at (7.5,-\i*1.32) {\f};
    \draw[ar] (q\i) -- (f\i);
  }
  \foreach \i [evaluate=\i as \j using int(\i+1)] in {1,...,5}{\draw[ar,teal] (n\i) -- (n\j);}
\end{tikzpicture}
\end{center}

\vspace{4pt}
{\headfont\bfseries\color{navy}Fonti e approfondimenti}
\begin{itemize}
  \item Lab originale: \emph{kathara lab -- basic IPv4 configuration, ping, traceroute and arp}, v3.0, Computer Networks Research Group, Università Roma Tre (\texttt{www.kathara.org}).
  \item Pagine di manuale di Kathará: \texttt{kathara(1)}, \texttt{kathara-lstart(1)}, \texttt{kathara-lconfig(1)}, \texttt{kathara-lab.conf(5)}, \dots{} su \texttt{www.kathara.org/man-pages}.
  \item Repository e documentazione: \texttt{github.com/KatharaFramework/Kathara}.
  \item ARP: RFC 826. ICMP: RFC 792.
\end{itemize}
